\section{A Separate Voice Layer}
\sectionkicker{VOICE FRONTIER}
\begin{wideflow}
{\Large\bfseries GPT-Live-1、独立した音声レイヤーとしてAPIに登場}\par
\smallskip
{\color{SurveyMuted}9月10日のGPT-Live-1を、Astraに同梱される機能ではなく、バックエンドのモデルに処理を委ねる独立した音声レイヤーとして描く。}\par
\end{wideflow}

9月10日、OpenAIはAPIでのGPT-Live-1を提供した\cite{gptlive}。聞き取りと発話を単一モデルで担う全二重の音声レイヤーで、深い推論やツール呼び出しはバックエンドのテキストモデルに委ねる構成である。バックエンドにはGPT-6 Astraのような自社モデルもサードパーティのモデルも置けるとされ、Astraのパッケージの一部として扱わない。割り込みへの応答や無音・雑音の処理、長時間セッションの維持、電話経由の利用がうたわれる。開発者はシステムプロンプトでトーンや速度、話し方を整え、音声認識のトランスクリプトやキーワードバイアスを使えるとされる。対応言語・方言や電話事業者の範囲は資料だけでは輪郭が定まらず、本号では言い切らない。

測定として示されるのは、Full Duplex BenchでGPT-Realtime-2.1を30ポイント上回ること、Astraの中程度設定と組み合わせてTau3で首位に立つこと、12を超える音声バリエーション、初期ユーザーの声として思考中の割り込み減少などである\cite{gptlive}。いずれもベンダー側のまとめであり、ベンチマークの手法やサンプリング、比較条件は今回たどっていない。料金は音声レイヤーとして1分あたり0.05ドルと示され、バックエンドのモデルやエージェントハーネスと組み合わせて使う構成である。カスタム音声は営業経由とされる。ユーザー企業の言葉はベンダー側が選んだ紹介として扱い、導入実績の実証とは切り分ける。

\begin{claimboundary}[音声レイヤーの境界]
ベンチマークの詳細はたどっていない。言語や電話の対応範囲は資料だけでは定まらない。ユーザー企業の言葉は選ばれた紹介である。
\end{claimboundary}
